\section{评测}

\definecolor{glm_bg}{RGB}{235, 245, 255}
\definecolor{glm_text}{RGB}{0, 110, 220}
\definecolor{header_gray}{RGB}{242, 242, 242}
\definecolor{section_gray}{RGB}{230, 230, 230}

\subsection{预训练评测}\label{sec:evals}

\begingroup
\renewcommand{\arraystretch}{1.1}

\begin{table*}[h]
    \centering
    \footnotesize
    \setlength{\tabcolsep}{2pt}
    \begin{tabular}{l c | c c c c c c}
        \toprule
        \rowcolor{header_gray}
			extbf{基准} & \textbf{示例数} &
        \makecell[c]{\textbf{\ourmodel}\\\textbf{Base}} &
        \makecell[c]{\textbf{MiMo-V2 Flash}\\\textbf{Base}} &
        \makecell[c]{\textbf{GLM-4.5}\\\textbf{Base}} &
        \makecell[c]{\textbf{DeepSeek V3.1}\\\textbf{Base}} &
        \makecell[c]{\textbf{DeepSeek V3.2}\\\textbf{Exp Base}} &
        \makecell[c]{\textbf{Kimi-K2}\\\textbf{Base}} \\
        \midrule
		激活参数量 & - & \glmcell{11B} & 15B & 32B & 37B & 37B & 32B \\
		总参数量 & - & \glmcell{196B} & 309B & 355B & 671B & 671B & 1043B \\
	        \midrule
			\rowcolor{section_gray}\multicolumn{8}{l}{\textsc{\textbf{通用}}} \\
	        BBH & 3-shot & \glmcell{88.2} & 88.5 & 86.2 & 88.2$\dagger$ & 88.7$\dagger$  & \best{88.7}  \\
	        MMLU & 5-shot & \glmcell{85.8} & 86.7 & 86.1 & 87.4$\dagger$ & \best{87.8}$\dagger$  & 87.8  \\
	        MMLU-Redux & 5-shot & \glmcell{89.2} & \best{90.6} & \na & 90.0$\dagger$ & 90.4$\dagger$  & 90.2  \\
	        MMLU-Pro & 5-shot & \glmcell{62.3} & \best{73.2} & \na & 58.8$\dagger$   & 62.1$\dagger$  & 69.2\\
	        HellaSwag & 10-shot & \glmcell{90.2} & 88.5 & 87.1 & 89.2$\dagger$   & 89.4$\dagger$  & \best{94.6}  \\
	        WinoGrande & 5-shot & \glmcell{79.1} & 83.8 & \na & \best{85.9}$\dagger$ & 85.6$\dagger$  & 85.3  \\
	        GPQA & 5-shot & \glmcell{41.7} & 43.5* & 33.5* & 43.1* & 37.3* & 43.1* \\
	        SuperGPQA & 5-shot & \glmcell{41.0} & 41.1 & - & 42.3$\dagger$  & 43.6$\dagger$  & \best{44.7} \\
	        SimpleQA & 5-shot & \glmcell{31.6} & 20.6 & 30.0 & 26.3$\dagger$ & 27.0$\dagger$  & \best{35.3}  \\
			\rowcolor{section_gray}\multicolumn{8}{l}{\textsc{\textbf{数学}}} \\
	        GSM8K & 8-shot & \glmcell{88.2} & \best{92.3} & 87.6 & 91.4$\dagger$ & 91.1$\dagger$  & 92.1  \\
	        MATH & 4-shot & \glmcell{66.8} & \best{71.0} & 62.6 & 62.6$\dagger$  & 62.5$\dagger$  & 70.2 \\
			\rowcolor{section_gray}\multicolumn{8}{l}{\textsc{\textbf{代码}}} \\
	        HumanEval & 3-shot & \glmcell{81.1} & 77.4* & 79.8* & 72.5* & 67.7* & \best{84.8}* \\
	        MBPP & 3-shot & \glmcell{79.4} & 81.0* & 81.6* & 74.6* & 75.6* & \best{89.0}* \\
	        HumanEval+ & 0-shot & \glmcell{72.0} & 70.7 & \na & 64.6$\dagger$  & 67.7$\dagger$  & \na \\
	        MBPP+ & 0-shot & \glmcell{70.6} & 71.4 & \na & \best{72.2}$\dagger$ & 69.8$\dagger$  & \na \\
	        MultiPL-E HumanEval & 0-shot & \glmcell{67.7} & 59.5 & \na & 45.9$\dagger$ & 45.7$\dagger$  & 60.5 \\
	        MultiPL-E MBPP & 0-shot & \glmcell{58.0} & 56.7 & \na & 52.5$\dagger$ & 50.6$\dagger$  & \best{58.8}  \\
			\rowcolor{section_gray}\multicolumn{8}{l}{\textsc{\textbf{中文}}} \\
	        C-EVAL & 5-shot & \glmcell{89.6} & 87.9 & 86.9 & 90.0$\dagger$ & 91.0$\dagger$  & \best{92.5}  \\
	        CMMLU & 5-shot & \glmcell{88.9} & 87.4 & \na & 88.8$\dagger$  & 88.9$\dagger$  & \best{90.9} \\
	        C-SimpleQA & 5-shot & \glmcell{63.2} & 61.5 & 70.1 & 70.9$\dagger$ & 68.0$\dagger$  & \best{77.6}  \\
	        \bottomrule
	    \end{tabular}
			\caption{预训练评测结果。* 表示原始分数不可用；为公平比较，我们在与 \ourmodel\ 相同的测试条件下复现并报告结果。$\dagger$ 表示 DeepSeek 分数引用自 MiMo-V2-Flash 报告~\cite{xiaomi2025mimo}。}
	    \label{tab:pretrain_eval_official}
\end{table*}

\endgroup


\paragraph{评测设置。} 我们在一系列覆盖不同能力的基准上评估 \ourmodel：
(1) 通用语言理解与推理，包括 BBH~\cite{suzgun2022challengingbbh}、MMLU~\cite{hendrycks2020mmlu}、MMLU-Redux~\cite{gema2024donewithmmlu}、MMLU-Pro~\cite{wang2024mmlupro}、HellaSwag~\cite{zellers2019hellaswag}、WinoGrande~\cite{sakaguchi2019winogrande}、GPQA~\cite{rein2023gpqa}、SuperGPQA~\cite{du2025supergpqa} 与 SimpleQA~\cite{openai2024simpleqa}。
(2) 数学推理，包括 GSM8K~\cite{cobbe2021gsm8k} 与 MATH~\cite{hendrycks2021math}。
(3) 编码能力，包括 HumanEval~\cite{chen2021evaluatinglargelanguagemodels}、MBPP~\cite{austin2021programsynthesislargelanguage}、HumanEval+、MBPP+~\cite{liu2023evalplus} 与 MultiPL-E~\cite{cassano2022multipl}。
(4) 中文理解，包括 C-Eval~\cite{huang2023ceval}、CMMLU~\cite{li2023cmmlu} 与 C-SimpleQA~\cite{he2024chinesesimpleqa}。

\paragraph{评测结果。}
表~\ref{tab:pretrain_eval_official} 总结了 \ourmodel 在通用推理、数学、代码与中文基准上的预训练评测。
尽管仅激活 11B 参数（总参数 196B），\ourmodel 仍与显著更大的稀疏基线（激活 15--37B；总参数 309--1043B）保持广泛竞争力，体现了强劲的准确率—效率权衡。
在核心通用基准上，\ourmodel 在 BBH 上达到 88.2（距最佳仅 0.5），在 MMLU 上达到 85.8。
值得注意的是，\ourmodel 在 SimpleQA 上达到 31.6，超过 DeepSeek-V3.2-Exp Base（27.0），且总参数仅 196B，相比 671B（约 $\sim$3.4$\times$），显示出更高的参数效率密度。
\ourmodel 还展现出强大的编码能力：HumanEval 81.1、MultiPL-E HumanEval 67.7、MultiPL-E MBPP 58.0。总体而言，这些结果表明 \ourmodel 在单位激活计算下具备高性能，为下游推理与智能体后训练提供坚实基础。

\subsection{后训练评测}
我们在代表性基准上评估 Step 3.5 Flash，包括推理导向的 HLE（文本子集）~\cite{phan2025hle}、MMLU-Pro~\cite{wang2024mmlupro}、GPQA-Diamond~\cite{rein2023gpqa}、AIME2025~\cite{10.5555/3524938.3525416}、HMMT~\cite{hmmt25}、IMO-AnswerBench~\cite{luong2025towards}；
代码相关的 LiveCodeBench-v6（2024.08–2025.05）~\cite{jain2024livecodebench}、CF-Div2-Stepfun\footnote{https://huggingface.co/datasets/stepfun-ai/CF-Div2-Stepfun}、SWE-Bench Verified~\cite{swe_verified} 与 SWE-Bench Multilingual~\cite{yang2025swesmith}；
智能体系列 $\tau^2$-Bench~\cite{tau2bench2025}、Terminal-Bench 2.0~\cite{terminalbench2025}、GAIA~\cite{mialon2023gaiabenchmarkgeneralai}、BrowseComp~\cite{wei2025browsecomp}、xbench-DeepSearch~\cite{chen2025xbench}、BrowseComp-zh~\cite{zhou2025browsecompzhbenchmarkingwebbrowsing} 与 \textsc{ResearchRubrics}~\cite{sharma2025researchrubrics}；
通用相关的 ArenaHard v2~\cite{arenahard2024}、IFBench~\cite{pyatkin2025generalizing} 与 MultiChallenge~\cite{sirdeshmukh2025multichallengerealisticmultiturnconversation}；以及长上下文相关的 LongBench v2~\cite{bai2024longbenchv2}、MRCR~\cite{vodrahalli2024michelangelo}~\footnote{https://huggingface.co/datasets/openai/mrcr}、FRAMES~\cite{krishna2025fact} 与 RepoQA~\cite{liu2024repoqaevaluatinglongcontext}。

我们还采用 \textbf{并行协同推理（PaCoRe）}范式~\cite{hu2026pacorelearningscaletesttime}，在推理、通用与长上下文基准上研究 Step 3.5 Flash 的测试时扩展特性。
依托 Step 3.5 Flash 的极致推理效率，该方法通过并行推理轨迹并在多轮协调中综合其洞见，将推理能力与上下文限制解耦。
具体而言，我们采用多轮 PaCoRe 轨迹配置 $\vec{K} = [4,4,4,4]$，在多个基准上获得显著增益。

我们保持最大序列长度 256k，使用默认解码配置（温度与 top-p 均为 1.0）。
在原始 128k 位置嵌入之上，我们应用 YaRN~\cite{peng2023yarnefficientcontextwindow}，缩放因子为 2.0，并仅作用于全注意力层。
所有方法均报告 pass@1 准确率，基于每题多次独立生成的平均表现：AIME 2025、HMMT 2025 Feb. 与 HMMT 2025 Nov. 为 64 次；IMO-AnswerBench、LiveCodeBench、GPQA-Diamond 与 MultiChallenge 为 8 次；HLE 为 1 次，其余基准为 4 次。
更多细节见附录~\ref{sec:post_training_eval_appendix}。

\paragraph{评测结果。} 表~\ref{fig:post-training-main-table} 给出了 Step~3.5~Flash 在推理、代码智能体、通用智能体、长上下文理解与通用能力等基准上的全面对比。
尽管仅激活 11B 参数（总参数 196B），Step~3.5~Flash 仍在广泛任务上表现强劲，尤其在 AIME~2025、HMMT~2025~Feb.、HMMT~2025~Nov.、IMO-AnswerBench 与 LiveCodeBench-v6 等高强度推理基准上突出。
它持续优于参数更多的开源模型，并达到 GPT-5.2~xHigh 与 Gemini~3.0~Pro 等前沿模型的同等级表现。
值得注意的是，Step~3.5~Flash 在智能体评测上也取得了强结果，包括 SWE-Bench Verified、Terminal-Bench~2.0、BrowseComp（含上下文管理）、GAIA 与 $\tau^2$-Bench，凸显其稳健的工具使用与长程决策能力。

\begingroup
\renewcommand{\arraystretch}{1.1}
\newlength{\stepcolw}
\setlength{\stepcolw}{1.45cm}
\newcolumntype{S}{>{\centering\arraybackslash}m{\stepcolw}}

\newcommand{\stepcell}[1]{\cellcolor{glm_bg}\textcolor{glm_text}{\textbf{#1}}}
\newcommand{\glmtext}[1]{\textcolor{glm_text}{\textbf{#1}}}

\begin{table}[t!]
    \centering
    \resizebox{\linewidth}{!}{
    \begin{tabular}{l|SS ccccc|ccc}
        \toprule
        \rowcolor{header_gray}
        \textbf{基准} &
        \multicolumn{2}{|c}{
          \makecell[c]{
                \\[0.0em]
        {\cellcolor{header_gray}\bfseries\fontsize{12pt}{14pt}\selectfont Step 3.5 Flash}\\[0.2em]
            \scriptsize
            \begin{tabular}{@{}>{\centering\arraybackslash}m{\stepcolw}@{\hspace{10pt}}>{\centering\arraybackslash}m{\stepcolw}@{}}
           \textbf{原始} & \textbf{PaCoRe}
        \end{tabular}
          }
        } &
        \makecell[c]{\textbf{MiniMax}\\\textbf{M2.1}} &
        \makecell[c]{\textbf{MiMo V2}\\\textbf{Flash}} &
        \makecell[c]{\textbf{GLM}\\\textbf{4.7}} &
        \makecell[c]{\textbf{DeepSeek}\\\textbf{V3.2}} &
        \textbf{\textbf{Kimi K2.5}} &
        \makecell[c]{\textbf{Gemini}\\\textbf{3.0 Pro}} &
        \makecell[c]{\textbf{Claude}\\\textbf{Opus}\\\textbf{4.5}} &
        \makecell[c]{\textbf{GPT-5.2}\\\textbf{xHigh}} \\
        \midrule
        激活参数量 &
        \multicolumn{2}{|c}{\cellcolor{glm_bg}\textcolor{glm_text}{\textbf{11B}}} &
        10B & 15B & 32B & 37B & 32B & - & - & - \\
        总参数量 &
        \multicolumn{2}{|c}{\cellcolor{glm_bg}\textcolor{glm_text}{\textbf{196B}}} &
        230B & 309B & 355B & 671B & 1T & - & - & - \\
        \midrule
        \rowcolor{section_gray}
        \multicolumn{11}{l}{\textsc{\textbf{推理}}} \\
        AIME 2025                & \glmcell{97.3} & \glmcell{99.9}  & 83.0 & 95.1* & 95.7  & 93.1   & 96.1 & 95.0 & 92.8 & \best{100.0} \\
        HMMT 2025 Feb.           & \glmcell{98.4} & \glmcell{100.0} & 71.0* & 95.4* & 97.1 & 92.5   & 95.4 & 97.5$\dagger$ & 92.9$\dagger$ & \best{99.4}\\
        HMMT 2025 Nov.           & \glmcell{94.0} & \glmcell{97.8}  & 74.3* & 91.0* & 93.5 & 90.2   & 91.1 & 94.5$\dagger$ & 91.7* & \best{97.1}* \\
        IMO-AnswerBench           & \glmcell{85.4} & \glmcell{88.8}  & 60.4* & 80.9* & 82.0 & 78.3   & 81.8 & 83.3$\dagger$ & 84.0$\dagger$ & \best{86.3}$\dagger$ \\
        LiveCodeBench-v6         & \glmcell{86.4} & \glmcell{88.9}  & 75.4* & 81.6* & 84.9 & 83.3   & 85.0 & \best{90.7}$\dagger$ & 84.8$\dagger$ & 87.7$\dagger$ \\
        CF-Div2-Stepfun-cpp         & \glmcell{86.1} & \glmcell{93.3}  & 59.0* & 46.9* & 74.1* & 81.6* & 73.6* & 83.5* & 72.2* & - \\
        MMLU-Pro                 & \glmcell{84.4} & \glmcell{84.8}  & 88.0 & 84.9 & 84.3 & 85.0     & 87.1 & \best{90.1}$\dagger$ & 89.5$\dagger$ & 87.4$\dagger$ \\
        GPQA-Diamond             & \glmcell{83.5} & \glmcell{85.0}  &  83.0 & 84.1* & 85.7 & 82.4   & 87.6 & 91.9 & 87.0 & \best{92.4} \\
        $\text{HLE}_\text{text}$ & \glmcell{23.1} & \glmcell{27.9} & 22.2 & 22.1 & 24.8 & 25.1      & 31.5 & \best{37.7}$\dagger$ & 30.8$\dagger$ & 35.5$\dagger$ \\
        \midrule

        \rowcolor{section_gray}
        \multicolumn{11}{l}{\textsc{\textbf{代码智能体}}} \\

        SWE Verified       & \glmcell{74.4} & \glmcell{\na} & 74.0 & 73.4 & 73.8 & 73.1 & 76.8 & 76.2 & \best{80.9} & 80.0 \\
        SWE Multilingual   & \glmcell{67.4} & \glmcell{\na} & 72.5 & 71.7 & 66.7 & 70.2 & 73.0 & 65.0$\dagger$ & \best{77.5}$\dagger$ & 72.0$\dagger$ \\
        Terminal-Bench 2.0 & \glmcell{51.0} & \glmcell{\na} & 47.9 & 38.5 & 41.0 & 46.4 & 50.8 & 56.9$\dagger$ & \best{59.3}$\dagger$ & 54.0$\dagger$ \\
        \midrule

        \rowcolor{section_gray}
        \multicolumn{11}{l}{\textsc{\textbf{通用智能体}}} \\

        BrowseComp              & \glmcell{51.6} & \glmcell{\na} & 47.4  & 45.4 & 52.0 & 51.4     & \best{60.6} & 37.8$\dagger$ & 37.0$\dagger$ & \na \\
        BrowseComp (w. Ctx Manage)      & \glmcell{69.0} & \glmcell{\na} & 62.0  & 58.3 & 67.5 & 67.6     & \best{74.9} & 59.2$\dagger$ & 57.8$\dagger$ & 65.8 \\
        BrowseComp-ZH           & \glmcell{66.9} & \glmcell{\na} & 47.8* & 51.2* & 66.6 & 65.0    & 62.3* & 66.8* & 62.4* & \best{76.1*} \\
        GAIA                    & \glmcell{84.5} & \glmcell{\na} & 64.3* & 78.2* & 61.9* & 75.1*  & 75.9* & 76.6* & 76.1* & 83.5* \\
        xbench-DeepSearch-2505               & \glmcell{83.7} & \glmcell{\na} & 68.7* & 69.3* & 72.0* & 78.0*  & 76.7* & 78.3* & 77.0* & 83.0* \\
        xbench-DeepSearch-2510               & \glmcell{56.3} & \glmcell{\na} & 43.0* & 44.0* & 52.3* & 55.7*  & 40.0$\dagger$ & 57.7* & 59.3* & \best{67.0*}\\
        \textsc{ResearchRubrics}& \glmcell{65.3} & \glmcell{\na} & 60.2* & 54.3* & 62.0* & 55.8*  & 59.5* & 50.1* &  61.6* & 57.8* \\
        $\tau^2$-Bench          & \glmcell{88.2} & \glmcell{\na} & 86.6* & 84.1* & 87.4 & 85.2*   & 85.4* & 90.7 & \best{92.5} & 85.5* \\

        \midrule
        \rowcolor{section_gray}
        \multicolumn{11}{l}{\textsc{\textbf{通用}}} \\
        Arena-Hard-v2.0 & \glmcell{74.0} & \glmcell{93.1} & 63.1* & 68.2* & 73.1* & 66.0*                         & \best{85.8}* & 81.7$\dagger$ & 76.7$\dagger$ & 80.6$\dagger$ \\
        MultiChallenge  & \glmcell{55.7} & \glmcell{60.8} & 50.5* & 44.3* & 67.8* & 57.1*                         & \best{73.6}* & 71.8* & 65.8* & 71.9* \\
        IFBench         & \glmcell{67.4} & \glmcell{56.8} & 70.0 & 64.0$\dagger$  & 68.0$\dagger$ & 61.0$\dagger$ & 72.8* & 70.4$\dagger$ & 58.0$\dagger$ & \best{75.4}$\dagger$ \\

        \midrule
        \rowcolor{section_gray}
        \multicolumn{11}{l}{\textsc{\textbf{长上下文}}} \\
        LongBench v2       & \glmcell{57.5}  & \glmcell{62.0} & 53.9* & 60.6$\dagger$ & 59.1* & 58.4$\dagger$                 & 61.0 & \best{70.0}* & 67.8* & 62.4* \\
        MRCR-8needle       & \glmcell{28.8}  & \glmcell{26.3} & 20.0$\dagger$ & 19.9$\dagger$ & 25.4$\dagger$ & 27.2$\dagger$ & 36.5* & 73.0$\dagger$ & 54.0* & \best{88.2*} \\
        FRAMES-Oracle      & \glmcell{76.5} & \glmcell{77.2} & 76.5* & 78.0* & 75.1* & 80.1*                                  & 77.4* & 79.7* & 85.8* & \best{87.3*} \\
        RepoQA             & \glmcell{88.5} & \glmcell{88.7} & 88.2* & 91.2* & 89.5* & 91.9*                                  & 89.8* & 91.5* & \best{95.7*} & 93.8* \\
        \bottomrule
    \end{tabular}}
    \caption{
    Step 3.5 Flash 与闭源/开源模型的对比。* 表示原始分数不可用或低于我们复现结果；为公平比较，我们在与 Step 3.5 Flash 相同的测试条件下复现并报告结果。$\dagger$ 表示分数来源于非官方渠道，包括技术报告或独立评测平台。我们对 HLE 的评测仅覆盖文本子集。BrowseComp（w. Ctx Manage）表示启用上下文管理的 BrowseComp 评测。
    }
    \label{fig:post-training-main-table}
\end{table}
\begin{comment}
\begin{table}[t!]
    \centering
    \footnotesize % 缩小字号以适应页面
    \setlength{\tabcolsep}{1.3pt} % 调整列间距
    \begin{tabular}{l|c c| c c c c c | c c c}
        \toprule
        \rowcolor{header_gray}
        \textbf{Benchmark} & 
        \makecell[c]{\textbf{Step 3.5}\\\textbf{Flash}} & 
        \makecell[c]{\textbf{with}\\\textbf{PaCoRe}} & 
        \makecell[c]{\textbf{MiniMax}\\\textbf{M2.1}} & 
        \makecell[c]{\textbf{MiMo}\\\textbf{V2 Flash}} & 
        \makecell[c]{\textbf{GLM}\\\textbf{4.7}} &
        \makecell[c]{\textbf{DeepSeek}\\\textbf{V3.2}} & 
        \makecell[c]{\textbf{Kimi}\\\textbf{K2.5}} &
        \makecell[c]{\textbf{Gemini}\\\textbf{3.0 Pro}} & 
        \makecell[c]{\textbf{Claude}\\\textbf{Opus 4.5}} & 
        \makecell[c]{\textbf{GPT-5.2}\\\textbf{Thinking}\\\textbf{(xhigh)}} \\
        \midrule
    
        \#Activated params & \glmcell{11B} &  & 10B & 15B & 32B & 37B & 32B & - & - & - \\
        \#Total params     & \glmcell{196B}&  & 230B& 309B& 355B& 671B& 1T  & - & - & - \\
        \midrule
    
        \rowcolor{section_gray}
        \multicolumn{11}{l}{\textsc{\textbf{Reasoning}}} \\
        AIME 2025                & \glmcell{97.3} & 99.9  & 83.0 & 95.1* & 95.7 & 93.1 & 96.1 & 95.0 & 92.8 & \best{100.0} \\
        HMMT 2025 Feb.           & \glmcell{98.4} & 100.0 & 71.0* & 95.4* & 97.1 & 92.5 & 95.4 & 97.5$\dagger$ & 92.9$\dagger$ & \best{99.4} \\
        HMMT 2025 Nov.           & \glmcell{94.0} & 97.8  & 74.3* & 91.0* & 93.5 & 90.2 & 91.1 & 94.5$\dagger$ & 91.7* & \best{97.1}* \\
        IMOAnswerBench           & \glmcell{85.4} & 88.8  & 60.4* & 80.9* & 82.0 & 78.3 & 81.8 & 83.3$\dagger$ & 84.0$\dagger$ & \best{86.3}$\dagger$ \\
        LiveCodeBench-v6         & \glmcell{86.4} & 88.9  & 75.4* & 81.6* & 84.9 & 83.3 & 85.0 & \best{90.7}$\dagger$ & 84.8$\dagger$ & 87.7$\dagger$ \\
        CF-Div2-Step-cpp         & \glmcell{86.1} & 93.3  & 59.0* & 46.9* & 74.1* & 81.6* & 73.6* & 83.5* & 72.2* & - \\
        MMLU-Pro                 & \glmcell{84.4} & 84.8  & 88.0 & 84.9 & 84.3 & 85.0 & 87.1 & \best{90.1}$\dagger$ & 89.5$\dagger$ & 87.4$\dagger$ \\
        GPQA-Diamond             & \glmcell{83.5} & 85.0  & 83.0 & 84.1* & 85.7 & 82.4 & 87.6 & 91.9 & 87.0 & \best{92.4} \\
        $\text{HLE}_\text{text}$ & \glmcell{23.1} & 27.9 & 22.2 & 22.1 & 24.8 & 25.1 & 31.5 & \best{37.7}$\dagger$ & 30.8$\dagger$ & 35.5$\dagger$ \\
        \midrule
    
        \rowcolor{section_gray}
        \multicolumn{11}{l}{\textsc{\textbf{Code Agent}}} \\
        SWE Verified       & \glmcell{74.4} & - & 74.0 & 73.4 & 73.8 & 73.1 & 76.8 & 76.2 & \best{80.9} & 80.0 \\
        SWE Multilingual   & \glmcell{67.4} & - & 72.5 & 71.7 & 66.7 & 70.2 & 73.0 & 65.0$\dagger$ & \best{77.5}$\dagger$ & 72.0$\dagger$ \\
        Terminal Bench 2.0 & \glmcell{51.0} & - & 47.9 & 38.5 & 41.0 & 46.4 & 50.8 & 56.9$\dagger$ & \best{59.3}$\dagger$ & 54.0$\dagger$ \\
        \midrule
    
        \rowcolor{section_gray}
        \multicolumn{11}{l}{\textsc{\textbf{General Agent}}} \\
        BrowseComp         & \glmcell{51.6} & - & 47.4 & 45.4 & 52.0 & 51.4 & \best{60.6} & 37.8$\dagger$ & 37.0$\dagger$ & \na \\
        BrowseComp (w. C.) & \glmcell{69.0} & - & 62.0 & 58.3 & 67.5 & 67.6 & \best{74.9} & 59.2$\dagger$ & 57.8$\dagger$ & 65.8 \\
        BrowseComp-ZH      & \glmcell{66.9} & - & 47.8* & 51.2* & 66.6 & 65.0 & 62.3* & 66.8* & 62.4* & \na \\
        GAIA               & \glmcell{84.5} & - & 64.3* & 78.2* & 61.9* & 75.1* & 75.9* & 76.6* & 76.1* & \na \\
        xbench-DeepSearch-2505          & \glmcell{83.7} & - & 68.7* & 69.3* & 72.0* & 78.0* & 76.7* & 78.3* & 77.0* & \na \\
        xbench-DeepSearch-2510          & \glmcell{56.3} & - & 43.0* & 44.0* & 52.3* & 55.7* & 40.0$\dagger$ & 57.7* & \best{59.3*} & \na \\
        \textsc{ResearchRubrics}
                           & \glmcell{65.3} & - & 60.2* & 54.3* & 62.0* & 55.8* & 59.5* & \na & 61.6* & \na \\
        $\tau^2$-Bench     & \glmcell{88.2} & - & 86.6* & 84.1* & 87.4 & 85.2* & 85.4* & 90.7 & \best{92.5} & 85.5* \\
        \midrule
    
        \rowcolor{section_gray}
        \multicolumn{11}{l}{\textsc{\textbf{General}}} \\
        Arena-Hard-v2.0 & \glmcell{74.0} & 93.1 & 63.1* & 68.2* & 73.1* & 66.0* & \best{85.8}* & 81.7$\dagger$ & 76.7$\dagger$ & 80.6$\dagger$ \\
        MultiChallenge  & \glmcell{55.7} & 60.8 & 50.5* & 44.3* & 67.8* & 57.1* & \best{73.6}* & 71.8* & 65.8* & 71.9* \\
        IFBench         & \glmcell{67.4} & 56.8    & 70.0 & 64.0$\dagger$ & 68.0$\dagger$ & 61.0$\dagger$ & 72.8* & 70.4$\dagger$ & 58.0$\dagger$ & \best{75.4}$\dagger$ \\
        \midrule
    
        \rowcolor{section_gray}
        \multicolumn{11}{l}{\textsc{\textbf{Long Context}}} \\
        LongBench v2  & \glmcell{57.5} & 62.0 & 53.9* & 60.6$\dagger$ & 59.1* & 58.4$\dagger$ & 61.0 & \best{70.0}* & 67.8* & 62.4* \\
        MRCR-8needle  & \glmcell{28.8} & 26.3 & 20.0$\dagger$ & 19.9$\dagger$ & 25.4$\dagger$ & 27.2$\dagger$ & 36.5* & 73.0$\dagger$ & 54.0* & \best{88.2*} \\
        FRAMES-Oracle & \glmcell{76.5} & 77.2 & 76.5* & 78.0* & 75.1* & 80.1* & 77.4* & 79.7* & 85.8* & \best{87.3*} \\
        RepoQA        & \glmcell{88.5} & 88.7 & 88.2* & 91.2* & 89.5* & 91.9* & 89.8* & 91.5* & \best{95.7*} & 93.8* \\
        \bottomrule
    \end{tabular}
    \caption{Comparison between Step 3.5 Flash and closed/open models. * denotes cases where the original score was unavailable or inferior to our reproduced result; we therefore report results evaluated under the same test conditions as Step 3.5 Flash for fair comparison. $\dagger$ indicates scores quoted from non-official sources, including technical reports, or independent evaluation platforms. Our evaluation on HLE focuses on the text-only subset.}
    \label{fig:post-training-main-table}
\end{table}
\end{comment}

\begin{comment}
\begin{table}[t!]
    \centering
    \footnotesize % 缩小字号以适应页面
    \setlength{\tabcolsep}{1.1pt} % 调整列间距
    \begin{tabular}{l|c|c c c c c c | c c c}
        \toprule
        % --- 表头 ---
        \rowcolor{header_gray}
        \textbf{Benchmark} & 
        \makecell[c]{\textbf{Step 3.5 Flash}} & 
        
        \makecell[c]{\textbf{MiniMax}\\\textbf{M2.1}} & 
        \makecell[c]{\textbf{MiMo V2}\\\textbf{Flash}} & 
        \makecell[c]{\textbf{GLM}\\\textbf{4.7}} &
        \makecell[c]{\textbf{DeepSeek}\\\textbf{V3.2}} & 
        \makecell[c]{\textbf{Kimi K2}\\\textbf{Thinking}} & 
        % \makecell[c]{\textbf{Kimi K2.5}\\\textbf{}} & 
        \textbf{\textbf{Kimi K2.5}} &
        \makecell[c]{\textbf{Gemini}\\\textbf{3.0 Pro}} & 
        \makecell[c]{\textbf{Claude}\\\textbf{Opus}\\\textbf{4.5}} & 
        \makecell[c]{\textbf{GPT-5.2}\\\textbf{Thinking}\\\textbf{(xhigh)}} \\
        \midrule

        \#Activated params & \glmcell{11B}  &  10B & 15B & 32B & 37B & 32B & 32B & - & - & - \\
        \#Total params     & \glmcell{196B} &  230B & 309B & 355B & 671B & 1T & 1T & - & - & - \\
        \midrule
        % --- REASONING ---
        \rowcolor{section_gray}
        \multicolumn{11}{l}{\textsc{\textbf{Reasoning}}} \\
        AIME 2025                & \glmcell{97.3/99.9} & 83.0 & 95.1* & 95.7 & 93.1 & 94.5 & 96.1 & 95.0 & 92.8 & \best{100.0} \\
        HMMT 2025 Feb.           & \glmcell{98.4/100.0} & 71.0* & 95.4* & 97.1 & 92.5 & 89.4 & 95.4 & 97.5$\dagger$ & 92.9$\dagger$ & \best{99.4}\\
        HMMT 2025 Nov.           & \glmcell{94.0/97.8} & 74.3* & 91.0* & 93.5 & 90.2 & 89.2* & 91.1 & \best{94.5}$\dagger$ & - & \na \\
        IMOAnswerBench           & \glmcell{85.4/88.8} & 60.4* & 80.9* & 82.0 & 78.3 & 78.6 & 81.8 & 83.3$\dagger$ & 84.0$\dagger$ & \best{86.3}$\dagger$ \\
        LiveCodeBench-v6         & \glmcell{86.4/88.9} & \na & 81.6* & 84.9 & 83.3 & 83.1 & 85.0 & \best{90.7}$\dagger$ & 84.8$\dagger$ & 87.7$\dagger$ \\
        CF-Div2-Step-cpp         & \glmcell{86.1/93.3} & 59.0* & 46.9* & 74.1* & 81.6* & 67.9* & \na & 83.5* & - & - \\
        MMLU-Pro                 & \glmcell{84.4/84.8} & 88.0 & 84.9 & 84.3 & 85.0 & 84.6 & 87.1 & \best{90.1}$\dagger$ & 89.5$\dagger$ & 87.4$\dagger$ \\
        GPQA-Diamond             & \glmcell{83.5/85.0} &  83.0 & 84.1* & 85.7 & 82.4 & 84.5 & 87.6 & 91.9 & 87.0 & \best{92.4} \\
        $\text{HLE}_\text{text}$ & \glmcell{23.1/27.9} & 22.2 & 22.1 & 24.8 & 25.1 & 23.9 & 31.5 & \best{37.7}$\dagger$ & 30.8$\dagger$ & 35.5$\dagger$ \\
        \midrule

        % --- CODE AGENT ---
        \rowcolor{section_gray}
        \multicolumn{11}{l}{\textsc{\textbf{Code Agent}}} \\
        
        SWE Verified       & \glmcell{74.4} &  74.0 & 73.4 & 73.8 & 73.1 & 71.3 & 76.8 & 76.2 & \best{80.9} & 80.0 \\
        SWE Multilingual   & \glmcell{67.4} &  72.5 & 71.7 & 66.7 & 70.2 & 61.1 & 73.0 & 65.0 & \best{77.5} & 72.0$\dagger$ \\
        Terminal Bench 2.0 & \glmcell{51.0} &  47.9 & 38.5 & 41.0 & 46.4 & 35.7 & 50.8 & 56.9 & \best{59.3} & 54.0 \\
        \midrule

        % --- GENERAL AGENT ---
        \rowcolor{section_gray}
        \multicolumn{11}{l}{\textsc{\textbf{General Agent}}} \\
        
        BrowseComp              & \glmcell{51.6} &  47.4 & 45.4 & 52.0 & 51.4 & 41.5$\dagger$ & \best{60.6} & 37.8 & 37.0 & \na \\
        BrowseComp (w. C.)      & \glmcell{69.0} &  62.0 & 58.3 & 67.5 & 67.6 & 60.2 & 74.9 & 59.2 & 57.8 & \best{85.8} \\
        BrowseComp-ZH           & \glmcell{66.9} &  47.8* & 51.2* & 66.6 & 65.0 & 62.3 & 62.3* & \best{66.8*} & \na & \na \\
        GAIA                    & \glmcell{84.5} &  64.3* & \best{78.2*} & 61.9* & 75.1* & 75.6* & 75.9* & 76.6* & \na & \na \\
        xbench-DeepSearch-2505               & \glmcell{83.7} &  68.7* & 69.3* & 72.0* & 78.0* & 76.0* & 76.7* & \best{78.3*} & 77.0* & \na \\
        xbench-DeepSearch-2510               & \glmcell{56.3} &  43.0* & 44.0* & 52.3* & 55.7* & \na & 40.0$\dagger$ & 57.7* & \best{59.3*} & \na\\
        \textsc{ResearchRubrics}& \glmcell{65.3} &  60.2* & 54.3* & \best{62.0}* & 55.8* & 56.2* & 59.5* & \na &  \na & \na \\
        $\tau^2$-Bench          & \glmcell{88.2} &  80.2* & 84.1* & 87.4 & 84.5* & 74.3$\dagger$ & \na & 90.7 & \best{92.5} & \na \\
        $\tau^2$-Bench Telecom  & \glmcell{97.3} &  87.0 & 95.3 & 96.5* & 99.1* & 93.0$\dagger$ & 95.9$\dagger$ & 98.0$\dagger$ & 98.2 & \best{98.7} \\

        \midrule
        % --- GENERAL ---
        \rowcolor{section_gray}
        \multicolumn{11}{l}{\textsc{\textbf{General}}} \\
        Arena-Hard-v2.0 & \glmcell{74.0 / 93.1} & 63.1* & 68.2* & 73.1* & 66.0* & 75.7* & \na & \best{81.7}$\dagger$ & 76.7$\dagger$ & 80.6$\dagger$ \\
        MultiChallenge  & \glmcell{55.7 / 60.8} & 50.5* & 44.3* & 67.8* & 57.1* & 66.4* & \na & \best{71.8}* & 54.2$\dagger$ & 57.9$\dagger$ \\
        IFBench         & \glmcell{67.4}   & 70.0 & 64.0$\dagger$ & 68.0$\dagger$ & 61.0$\dagger$ & 68.0$\dagger$ & \na & 70.4$\dagger$ & 58.0$\dagger$ & \best{75.4}$\dagger$ \\

        \midrule
        % --- LONG CONTEXT ---
        \rowcolor{section_gray}
        \multicolumn{11}{l}{\textsc{\textbf{Long Context}}} \\
        LongBench-V2       & \glmcell{57.5}  & 53.9* & 60.6$\dagger$ & 59.1* & 58.4$\dagger$ & 48.1$\dagger$ & 61.0 & - & - & - \\
        MRCR-8needle(auc)  & \glmcell{28.8}  & 20.0$\dagger$ & 19.9$\dagger$ & 25.4$\dagger$ & 27.2$\dagger$ & 24.0* & \na & 73.0* & - & - \\
        FRAMES-ORACLE      & \glmcell{76.5}* & 76.4* & 78.0* & 75.1* & 80.1* & 76.9* & \na & 79.7* & - & - \\
        REPOQA             & \glmcell{88.5}* & 88.1* & 91.2* & 89.5* & 91.9* & 89.2* & \na & 91.5* & - & - \\
        \bottomrule
    \end{tabular}
    \caption{Comparison between Step 3.5 Flash and closed/open models. “*” denotes cases where the original score was unavailable or inferior to our reproduced result; we therefore report results evaluated under the same test conditions as Step 3.5 Flash for fair comparison. $\dagger$ indicates scores quoted from non-official sources, including technical reports, or independent evaluation platforms. Our evaluation on HLE focuses on the text-only subset.}
    \label{fig:post-training-main-table-script}
\end{table}
\end{comment}


\endgroup
